\documentclass[12pt]{article}
\usepackage{amsmath, amssymb}
\usepackage{geometry}
\geometry{margin=1in}

\title{Lecture 26: Stochastic Processes (Part 1)}
\author{}
\date{}

\begin{document}

\maketitle

\section*{Strict Sense Stationary (SSS), Wide Sense Stationary (WSS), and Ergodic Process}

\section{Basics of Stochastic Processes (Revision)}

\subsection{Collection of Waveforms}

\textbf{Definition (Stochastic Process)}  

A stochastic process is defined as:
\[
X(\zeta, t)
\]

\textbf{Logical Interpretation}  

An entire collection of waveforms or functions over time dependent on the outcome of a random experiment.

The definition emphasizes that:
\begin{itemize}
    \item The process is not a single function but a collection of possible functions (waveforms).
    \item Each waveform corresponds to an outcome of a random experiment.
    \item Therefore, the output is inherently random.
\end{itemize}

\textbf{Key Observation}  

The random output heavily depends on the experiment.

\textbf{Example (From Slides)}  

Tossing a coin:
\begin{itemize}
    \item Outcomes: Heads or Tails
    \item Each outcome produces a specific deterministic waveform
\end{itemize}

Hence:
\[
\text{Random experiment} \rightarrow \text{deterministic realization}
\]

\subsection{Realizations and Random Variables}

\textbf{Definition (Realization)}  

A specific measured waveform from the experiment is called a realization of the process.

\textbf{Fixed-in-Time Representation}

\textbf{Step 1: Fix Time}  

Consider a specific time $t = t_1$

\textbf{Step 2: Evaluate Process}  

The stochastic process reduces to a single random variable:
\[
X(\zeta, t_1) = X(t_1)
\]

\textbf{Step 3: Extend to Multiple Times}  

At multiple time instants $(t_1, t_2)$:
\[
X(t_1), X(t_2)
\]

These are multiple dependent random variables.

\textbf{Key Conclusion}  

The dependency between these random variables determines:

\section{Characterization of Stochastic Process}

\subsection{Probability Density Functions (PDF)}

\textbf{1st Order Characterization}  

Describes the random variable at a single time $t_1$:
\[
f_X(x, t_1), \quad F_X(x, t_1)
\]

\textbf{Reasoning}  

Since fixing time gives a random variable:
\[
\rightarrow \text{Its behavior is captured using PDF/CDF}
\]

\textbf{2nd Order Characterization}  

Describes joint probability at two time instants $(t_1, t_2)$:
\[
f_{X_1, X_2}(x_1, x_2, t_1, t_2)
\]

\textbf{Reasoning}  

Since multiple time evaluations give multiple dependent random variables:
\[
\rightarrow \text{Their joint distribution is required}
\]

\textbf{nth Order Characterization}  

Complete statistical description:
\[
f_{X_1, \dots, X_n}(x_1, \dots, x_n, t_1, \dots, t_n)
\]

\textbf{Conclusion}  

Increasing order $\rightarrow$ more complete description of the process

\section{Statistical Moments}

The fundamental nature of the stochastic process

\subsection{Mean of Stochastic Process (First Order Moment)}

\textbf{Definition}  

The mean of a stochastic process is:
\[
\mu_X(t) = E[X(t)]
\]

\textbf{Step-by-Step Computation}  

Using first-order PDF:
\[
\mu_X(t) = \int_{-\infty}^{\infty} x \cdot f_X(x, t)\, dx
\]

\textbf{Observation}  

The mean is generally a function of time:
\[
\mu_X(t) \neq \text{constant}
\]

\textbf{Reasoning}  

Since the distribution can vary with time:
\[
\rightarrow \text{The expected value also varies with time}
\]

\subsection{Autocorrelation Function}

\textbf{Purpose}  

Measures dependence between values at different time instants

\textbf{Definition}  
\[
R_{XX}(t_1, t_2) = E[X(t_1)X(t_2)]
\]

\textbf{Interpretation}  

Captures relationship between:
\[
X(t_1) \quad \text{and} \quad X(t_2)
\]

\textbf{Logical Development}
\begin{enumerate}
    \item Multiple time evaluations $\rightarrow$ multiple random variables
    \item These variables are dependent
    \item Their dependency is quantified by autocorrelation
\end{enumerate}

\section{Conceptual Flow of the Lecture}

The lecture builds logically as follows:
\begin{enumerate}
    \item Start with definition  

    Stochastic process = collection of waveforms

    \item Fix time  

    Leads to random variables

    \item Multiple time instants  

    Leads to dependent random variables

    \item Characterization  

    Use PDFs (1st, 2nd, nth order)

    \item Statistical summaries  

    Mean $\rightarrow$ first-order behavior  

    Autocorrelation $\rightarrow$ dependency structure
\end{enumerate}

\section{Summary for Exam Revision}

\begin{itemize}
    \item A stochastic process is a collection of time functions indexed by random outcomes
    \item Fixing time $\rightarrow$ random variable
    \item Multiple times $\rightarrow$ dependent random variables
    \item Full description requires nth-order PDFs
\end{itemize}

\textbf{Mean:}
\[
\mu_X(t) = E[X(t)]
\]

\textbf{Autocorrelation:}
\[
R_{XX}(t_1, t_2) = E[X(t_1)X(t_2)]
\]

\end{document}